% ============================================================
%  Section 6 — Discussion  (~1.5 pages)
% ============================================================
\section{Discussion}

\subsection{Framework Selection Guidance}

For NISQ practitioners selecting a framework for oracle, error
correction, or search workloads: Qiskit and Cirq are functionally
interchangeable in circuit-quality terms.
A 0.01 difference in $S_{\text{overall}}$ (58.80 vs.\ 58.79) lies below
any meaningful threshold; selection between them should be based on
ecosystem fit (IBM hardware $\to$ Qiskit; Google hardware $\to$ Cirq)
or team familiarity, not performance.
Both produce constant-depth circuits for oracle algorithms, maintain
statistical equivalence through 8 qubits, and compile with stable,
predictable latency.

PennyLane's advantages lie outside this study's scope.
Its automatic differentiation engine and device-agnostic \texttt{QNode}
abstraction make it the natural choice for variational algorithm research
(VQE, QAOA, QML), where gradient computation across hybrid
quantum-classical loops is the primary concern.
Practitioners using PennyLane for VQE and QML workloads at small qubit
counts ($\leq$2--3 qubits) will observe no equivalence penalty.
The trade-off---larger compiled circuits and JSD$=1.000$ failures at
scale---becomes operationally significant only when the algorithm
requires 4+ qubits and output distributions must be compared with those
from Qiskit or Cirq implementations.

\subsection{The Decomposition Incompatibility Problem}

The JSD$=1.000$ failures reported in Section~5 are not compilation
errors.
PennyLane produces syntactically valid, executable OpenQASM 3.0 circuits
that QSim runs without errors.
The issue is semantic: PennyLane's full decomposition of multi-qubit
oracle gates into primitive two-qubit gates changes the computational
basis in which results are measured at 4+ qubits.
Where Qiskit and Cirq implement a multi-controlled phase operation as a
single structured gate whose output matches the expected oracle behavior,
PennyLane expands the same gate into a sequence of elementary rotations
that, while equivalent in principle, produces a different effective
mapping between input and output basis states when composed with the
algorithm's measurement layer.

This is an underdocumented behavior with serious practical implications.
A practitioner who ports an algorithm from Qiskit to PennyLane at 5
qubits and compares output distributions will observe completely
non-overlapping histograms (JSD$=1.000$) and may incorrectly conclude
their implementation is wrong.
The correct diagnosis is a framework-induced decomposition
incompatibility that is invisible without a common backend.

\subsection{The Unified Pipeline as a Methodological Contribution}

The JSD$=1.000$ findings are discoverable only because all circuits
execute on the same QSim backend.
Had this study used PennyLane's \texttt{default.qubit} alongside
Qiskit's Aer, any divergence would be confounded by the two simulators'
different state-update algorithms, shot-sampling implementations, and
floating-point precision policies.
The QCanvas pipeline removes this confound entirely: a JSD$=1.000$
result in our setup cannot be attributed to the backend; it must be
attributed to the framework's gate-level choices.

\subsection{Hardware Topology and Connectivity Constraints}
\label{sec:hw-topology}

All experiments in this study assume ideal all-to-all qubit connectivity,
corresponding to statevector simulation with no routing overhead.
On real quantum hardware, qubits are arranged in restricted topologies
(e.g., IBM heavy-hex graphs, Google Sycamore grids) and two-qubit gates
can only be applied to physically adjacent qubits.

When a target coupling map is enforced, the transpiler must insert SWAP
gates to route logical qubits to compatible physical qubit pairs;
this routing step is the most computationally expensive and
quantum-error-prone component of the transpilation pipeline.
Comparing Qiskit and Cirq without a coupling map therefore ignores the
primary domain where their transpilers differ most substantially:
Qiskit's SABRE-based router~\cite{li2019tackling} and Cirq's
\texttt{LineSwapRouter} employ fundamentally different heuristic
strategies that produce very different SWAP-overhead profiles depending on
circuit width and hardware topology.
Future work should replicate this benchmark suite on a hardware-constrained
topology (e.g., a 27-qubit heavy-hex map) to characterise the routing
overhead and SWAP insertion penalty for each framework, which we expect
would significantly change the relative gate counts and circuit depths
observed in this study.

\subsection{Limitations}

The following limitations bound the generalisability of these results:

\begin{itemize}
  \item \textbf{Classical simulation only.} All results are from
        statevector simulation.
        Physical hardware validation on NISQ devices---where crosstalk,
        leakage, and device-specific noise profiles dominate---remains
        future work.
        No hardware implications are claimed beyond what the depolarizing
        noise model supports.
  \item \textbf{Simulator mode.} QSim operates in statevector mode;
        density matrix simulation and stabilizer-circuit simulation are
        not evaluated.
  \item \textbf{No hardware coupling map.} All benchmarks assume ideal
        all-to-all connectivity.
        Enforcing a realistic hardware topology would activate SWAP-gate
        insertion and routing heuristics---the most performance-sensitive
        part of any transpiler---which are not evaluated here.
        Gate counts and depths reported in Section~5 therefore represent
        lower bounds on the overhead incurred on physical devices
        (see Section~\ref{sec:hw-topology} for further discussion).
  \item \textbf{PennyLane version pinned at 0.35.1.} Results reflect the
        PennyLane release pinned in the benchmark environment
        (Table~\ref{tab:versions}).
        PennyLane's decomposition strategy may change in future releases.
  \item \textbf{Quantum Volume is structure-implied.} QV is computed from
        circuit structure under a depolarizing model, not from
        heavy-output probability on physical hardware.
        It does not capture real noise profiles (crosstalk, leakage).
  \item \textbf{Algorithm scope.} The 12-algorithm suite is
        representative but not exhaustive.
        Deeper QAOA layers, fault-tolerant codes, and quantum simulation
        circuits are out of scope.
  \item \textbf{Grover scaling extended to 15 qubits (compilation only).}
        To strengthen the power-law extrapolation beyond 5 simulation
        data points, gate-count measurements for Grover's Algorithm were
        extended to 15 qubits via compilation-only runs.
        Statevector simulation was not performed for $n \geq 7$ qubits
        due to memory constraints ($2^{15}$ amplitudes $\approx$256~GB).
\end{itemize}
